Rapid, all-weather mapping of snow avalanche debris across rugged mountain terrain is critical for alpine civil protection, emergency response, and search-and-rescue operations. Synthetic Aperture Radar (SAR) sensors such as Sentinel-1 provide cloud-penetrating day-and-night winter monitoring, but bitemporal change detection in steep terrain suffers from foreshortening, layover, and radar shadow. We propose \textbf{TopoRadar-Net}, coupling bitemporal Sentinel-1 C-band SAR with local radar-topographic geometry through: (1) an analytical characterization of shift-invariant linear-convolution limits under non-stationary terrain projection; (2) multi-scale Radar-Topographic Cross-Attention Blocks; (3) a heuristic slope-prior loss audited for positive-label conflict; and (4) calibrated component- and corridor-level operational proxies. On AvalCD (7 events, 4 mountain systems, 993 polygons), TopoRadar-Net reaches $63.79\% \pm 1.88\%$ pooled macro F1, compared with Attention U-Net ($63.41\%$), SiamUNet-conc ($63.13\%$), ResU-Net ($62.92\%$), and SiamUNet-diff ($54.87\%$). Across three seeds, only the margin over SiamUNet-diff is significant ($+8.92\%$, $p=0.021$); the Attention U-Net difference is $+0.38\%$ ($p=0.869$). Seed-42 backslope differences do not persist as stable three-seed effects (Troms{\o}: $-0.34\% \pm 6.28\%$; Pamir: $+1.20\% \pm 7.73\%$). In Pamir, the false-alarm footprint is $30.0\%$ lower than Attention U-Net ($72.9\text{ ha}$ nominal-grid equivalent; $55.6\text{ ha}$ affine physical difference). The slope-prior mask overlaps $18.47\%$ of training positives and lowers held-out positive recall inside that mask relative to No-GeoLoss. Validation-calibrated Pamir road-corridor alerts achieve $83.0\%$ precision and $62.5\%$ recall, but the single Troms{\o} road-intersecting avalanche is missed in all seeds. These results establish competitive pixel performance and explicit operational failure boundaries, not a component-isolated mechanism or deployment-ready system.
